\documentclass[11pt]{article}
\usepackage{amsmath,amssymb,amsthm}
\newtheorem{lemma}{Lemma}
\newtheorem{corollary}{Corollary}
\newcommand{\Q}{\mathbb Q}
\newcommand{\Z}{\mathbb Z}
\begin{document}
\title{An exact product-height identity for odd elliptic indices}
\author{}\date{}\maketitle
Put $K=\Q(\sqrt2)$, $d=3-\sqrt2$, $D=d^2$, and
\[
 E:y^2=x^3-Dx,\qquad
 P=\left(\frac{57-28\sqrt2}{8},\frac{-660+417\sqrt2}{32}\right).
\]
Write $h$ for the absolute logarithmic Weil height. We use the source's
facts that $P$ has infinite order and belongs to $2E(K)$.

\begin{lemma}\label{lem:oddproduct}
For any finite list of odd integers $n_1,\ldots,n_r$, one has
\[
 h\!\left(\prod_{i=1}^r x([n_i]P)\right)
 =\sum_{i=1}^r h(x([n_i]P)).
\]
\end{lemma}
\begin{proof}
Put $T_0=(0,0)$. At every good finite place $v$, let $\bar P$ be the
reduction of $P$ and let $m$ be its order. For a nonidentity point $Q$,
$v(x(Q))<0$ means that $Q$ reduces to the identity, whereas
$v(x(Q))>0$ means that it reduces to $T_0$. The latter has order two.
If $m$ is odd, no multiple of $\bar P$ can equal $T_0$, so
$v(x([n_i]P))\le0$ for all $i$. If $m$ is even, an odd multiple of
$\bar P$ cannot be the identity, so $v(x([n_i]P))\ge0$ for all $i$.
Thus at a good place all the valuations have the same weak sign.

Here is a formulation which also treats the place above two without
requiring a smooth reduction model. Let $E_1(K_v)$ be the formal
subgroup, consisting of the identity and the points with $v(x)<0$.
It is an open subgroup of finite index. The image of $T_0$ in
$E(K_v)/E_1(K_v)$ has order two. For $Q\ne O,T_0$, the addition formula is
\[
 x(Q+T_0)=-D/x(Q).
\]
At every place with $v(D)=0$, positive valuation of $x(Q)$ is therefore
equivalent to $Q+T_0\in E_1(K_v)$. Apply the preceding order-parity
argument in the finite quotient $E(K_v)/E_1(K_v)$. It again shows that
all $v(x([n_i]P))$ have the same weak sign.

The only remaining place is the prime $\mathfrak d=(3-\sqrt2)$ above
seven. The displayed integral Weierstrass model is minimal there,
since $v_{\mathfrak d}(-D)=2<4$. Its reduction is the cuspidal cubic
$y^2=x^3$ over $\mathbf F_7$. The point $P$ reduces to $(1,6)$:
\[
 \frac{57-28\cdot3}{8}\equiv1,\qquad
 \frac{-660+417\cdot3}{32}\equiv6\pmod7.
\]
This point is nonsingular. Hence all multiples of $P$ belong to the
subgroup $E_0(K_{\mathfrak d})$ of points with nonsingular reduction.
If $v_{\mathfrak d}(x(Q))>0$, the curve equation implies
$v_{\mathfrak d}(y(Q))>0$, so $Q$ reduces to the singular point $(0,0)$.
Consequently $v_{\mathfrak d}(x([n_i]P))\le0$ for every $i$.
This completes the finite-place argument.

At each real embedding, $x([n_i]P)$ is strictly positive because
$[n_i]P\in2E(K)$. Its ordinate is nonzero. The equation of $E$ then gives
$x([n_i]P)>d$ at that embedding. Both values $3\pm\sqrt2$ of $d$ exceed
one. Thus all the real absolute values of these abscissae exceed one.

At a finite place where all valuations have the same weak sign,
\[
 \log\max\!\left(1,\left|\prod_i x([n_i]P)\right|_v\right)
 =\sum_i\log\max(1,|x([n_i]P)|_v).
\]
The same identity holds at each real embedding by the preceding
paragraph. Summing with the usual normalized local weights proves
the result.
\end{proof}

\begin{corollary}[Three detector scales]\label{cor:threescales}
Suppose that the four root points $Q_1,\ldots,Q_4$ in the shared-point
construction have odd integer indices. Write
\[
 C=\prod_{i=1}^4x(Q_i)=C_0+C_1\sqrt2.
\]
The nine height-detector scales can be replaced by the three scales
\[
 1,\qquad C_0^2,\qquad(1+C_1^2)^2.
\]
Let $\kappa$ be the radical exponent in the chosen source height theorem ($24$ in the historical shared-point source and $19$ in the current refinement). The point-height bound with the original constant $7H$ remains valid.
The source's simultaneous-prime and pole-completeness construction
continues to apply to these scales once the root points are fixed.
\end{corollary}
\begin{proof}
The element $C$ is totally positive, so $C_0>0$. Thus the three scales
are positive nonzero rational squares, exactly as required in the
completeness argument.

Let $N\ne0$ be an integer divisible by the $k$th power of each of the
three contact radicals, and put $T=(H+\log3)|N|^{\kappa/k}$. The detector
height transfer gives
\[
 h(\lambda),\ h(C_0^2\lambda),\ h((1+C_1^2)^2\lambda)\le T.
\]
Therefore $h(C_0)\le T$ and $h(C_1)\le T/2$, by the same elementary
rational height identities used for the original fixed scales. Hence
\[
 h(C)\le h(C_0)+h(C_1)+\tfrac32\log2
 \le\tfrac32T+\tfrac32\log2
 <7H|N|^{\kappa/k}.
\]
Lemma~\ref{lem:oddproduct} now gives $h(x(Q_i))\le h(C)$ for every $i$.
Thus the original bound for each root point holds with the same
constant. In particular, if $B$ is the largest absolute value of the
four indices, the original lower bound yields
$B^2-6<7H|N|^{\kappa/k}$. No numerical constant or anchor enlargement is
required for this replacement itself.

The source's completeness proof uses the scales only as fixed positive
rational squares after the root points have been selected. Its contact
positions, prime avoidance, common slope, and pole witnesses therefore
apply verbatim to this shorter list.
\end{proof}

\paragraph{Scope.}
The oddness of all four indices is essential to this proof. This note
does not assert that the existing source already enforces it. The
construction must supply an odd-index version of its Lucas step and
point predicates. Subject to that separate change, the detector
replacement above is an unconditional exact height argument.
\end{document}
